\section{Code Forensics}
\label{sec:code_forensics}

This section reconstructs the scientific rationale embedded in the Grand Prize winning codebase~\cite{nader2024gp, landau2024gpplus}, explaining \emph{why} key design decisions were made and how they relate to the physics of the ink detection problem.

\subsection{Data Loading: Non-Obvious Complexity}

The data loader handles two critical issues not visible from the pipeline interface. First, layer images use inconsistent naming conventions across segments (\texttt{00.tif}, \texttt{0.tif}, \texttt{layer\_00.tif}) and formats (TIFF for Scroll~1, JPEG for Scroll~5); the loader resolves these through priority-based pattern matching. Second, eight Scroll~1 segments have inverted Z-ordering-layer 0 faces the scroll exterior rather than the interior-a fact not documented in the data files and discovered only by examining depth profiles. Both issues fail silently without explicit handling, producing degraded predictions with no error signal.

\subsection{Z-Axis Augmentation: Domain-Driven Robustness}

The most instructive design choice is the ``fourth augmentation'' (\texttt{fourth\_augment}): three augmentations applied along the depth axis rather than spatially. Z-flip ($p = 0.5$) trains the model on both Z-orderings, directly addressing the reversed-segment problem. Z-shift ($\Delta z \in [-2, +2]$, $p = 0.3$) simulates uncertainty in the surface estimation, since the segmentation tool's estimated surface position varies locally by 1--3 voxels across a segment. Layer dropout ($p = 0.1$) builds robustness to missing or corrupted depth data. Together, these augmentations build invariance to the exact failure modes that cause real-world segment failures-a deliberate engineering choice rather than generic regularization.

\subsection{Gaussian Tile Reassembly}

Naive uniform averaging of overlapping tile predictions produces visible grid artifacts because tile-edge predictions are less reliable than tile-center predictions (truncated receptive field). Gaussian-weighted reassembly assigns high weight to predictions near each tile center and near-zero weight at boundaries ($\sigma = P/6 \approx 10.67$), producing a smooth, artifact-free output map. The approach is equivalent to a kernel density estimator over the tile grid: denser tiling produces smoother output as more Gaussian ``bumps'' overlap.

\subsection{TimeSformer: The Video-to-Depth Analogy}

The TimeSformer~\cite{bertasius2021timesformer} was designed for video classification, where temporal attention captures motion patterns across frames. In the ink detection context, Z-layers play the role of video frames: each layer shows the same papyrus region at a different depth, and temporal attention learns which cross-layer relationships carry ink signal. This analogy directly motivates the architecture choice-ink detection requires distinguishing surface ink (a depth-localized signal) from through-papyrus artifacts (diffuse across layers), precisely the cross-frame discrimination that temporal attention was built for. The analogy has limits: unlike video frames, adjacent CT layers can show abrupt structural changes at material boundaries, which may explain why the simpler U-Net, imposing no sequential structure, can match or exceed the TimeSformer with only 5 input layers.

\subsection{Design Philosophy}

The codebase reveals a coherent philosophy: \emph{build robustness through augmentation, not architectural complexity}. Domain-specific challenges (reversed segments, surface uncertainty, data quality variation) are addressed through explicit data augmentation rather than specialized model design. This simplicity-first approach is directly reproducible and serves as the scientific baseline this work formalizes.

